\chapter{A library of chance mechanisms}
\label{ch:standard-distributions}
\chaptertag{Discrete distributions first; continuous models later}

\begin{question}
Which probability laws recur because the same mechanisms recur?
\end{question}

For a discrete variable, its distribution lists the possible values and
their probabilities.
The coin counts, dice, and hat counts of Chapter 4 already have
distributions. A \emph{named family} collects experiments whose probability
tables have the same form. Learn the experiment before memorizing the name.

\begin{bigidea}[title={A route through the next topics}]
Start with finite distributions: a uniform choice, one success indicator,
a count of independent successes, and sampling without replacement.
For each, identify the possible values, calculate their probabilities,
then find the mean and variance.
Geometric waiting times and Poisson counts come next as countably infinite
examples. The coupon collector then brings waiting times, expectation,
variance, and inclusion--exclusion together. The continuous section is
for later, after the calculus refresher.
\end{bigidea}

\section{A uniform choice from a finite list}
\label{sec:discrete-families}

Choose one of $1,\ldots,m$ uniformly, where $m\ge1$. Then
$\P(U=k)=1/m$ for $k=1,\ldots,m$, and zero for other values.
The usual sum and sum-of-squares identities give
\[
 \E[U]=\frac{m+1}{2},\qquad
 \E[U^2]=\frac{(m+1)(2m+1)}6,\qquad
 \Var(U)=\frac{m^2-1}{12}.
\]
For example, choosing a label uniformly from $1,\ldots,8$ gives
mean $9/2$ and variance $21/4$.
The case $m=1$ is constant, with variance zero.

A uniform distribution on any finite set $\{u_1,\ldots,u_m\}$ averages
those values with weights $1/m$. The formulas above are specifically for
the consecutive integers $1,\ldots,m$.
Uniform outcomes also need not produce uniform recorded values:
the sum of two fair dice is not uniform on $2,\ldots,12$.

\section{Bernoulli: one event, recorded as zero or one}
\label{sec:bernoulli}

If an event has probability $p$, its indicator $I$ has distribution
\[
 \P(I=1)=p,\qquad \P(I=0)=1-p,\qquad 0\le p\le1.
\]
We call this the \emph{Bernoulli distribution} with parameter $p$.
There are only two values, so directly
\[
 \E[I]=p,\qquad \E[I^2]=p,\qquad \Var(I)=p(1-p).
\]
The parameter $p$ is a probability, the value of $I$ is zero or one,
and its expectation can lie between these values.
At $p=0$ or $1$ the result is certain and the variance vanishes.

\begin{example}
Let $I$ record whether a fair die shows a six. Then $p=1/6$,
$\E[I]=1/6$, and $\Var(I)=5/36$.
A success may mean a correct answer, a defective item, or a birthday
match for a specified pair; a \emph{success} need not mean something desirable.
\end{example}

\section{Binomial: count successes in a fixed number of trials}
\label{sec:binomial}

Perform $n$ independent trials, each with success probability $p$,
and let $X$ count the successes. For $0<p<1$ and $k=0,\ldots,n$,
\[
 \P(X=k)=\binom nkp^k(1-p)^{n-k}.
\]
A specified success/failure string has probability $p^k(1-p)^{n-k}$;
choose its $k$ success positions in $\binom nk$ ways.
The binomial theorem verifies that the masses sum to one.
We write $X\sim\Bin(n,p)$.
If $p=0$, $X=0$; if $p=1$, $X=n$; if $n=0$, $X=0$.

Write $X=I_1+\cdots+I_n$. Linearity and the independent-variance rule give
\[
 \E[X]=np,\qquad \Var(X)=np(1-p).
\]
This calculation avoids summing $k^2$ against the binomial probabilities.

\begin{example}[Five trials]
For $X\sim\Bin(5,1/3)$,
\[
 \P(X=3)=\binom53(1/3)^3(2/3)^2=\frac{40}{243},
 \qquad \E[X]=\frac53,\qquad \Var(X)=\frac{10}{9}.
\]
For a threshold it can be easier to use a complement:
$\P(X\ge2)=1-\P(X=0)-\P(X=1)$.
\end{example}

\begin{warningbox}
A count is binomial only when its experiment satisfies the assumptions:
a fixed number of trials, two outcomes per trial, a common success
probability, and independence.
Recording the same toss ten times gives a count of either zero or ten,
not a binomial count of ten independent tosses.
Sampling from a finite population without replacement also requires
a different model.
\end{warningbox}

\section{Hypergeometric: sample without replacement}
\label{sec:hypergeometric}

A population contains $K$ marked objects among $N$ distinct objects.
Choose a subset of $m$ objects uniformly, and let $X$ count marked objects
in the sample. There are $\binom Nm$ equally likely samples.
To obtain exactly $k$ marked objects, choose those $k$ from the marked
group and the other $m-k$ from the unmarked group:
\[
 \P(X=k)=\frac{\binom Kk\binom{N-K}{m-k}}{\binom Nm}.
\]
Here $N\ge1$, $0\le K,m\le N$, and the possible values satisfy
$\max(0,m-(N-K))\le k\le\min(m,K)$.
The formula is normalized because these disjoint cases partition all
$m$-element samples. This is the \emph{hypergeometric distribution}.

\begin{example}[Two draws from five balls]
An urn contains three red and two blue balls. For a uniform sample of
two balls without replacement, the number $X$ of red balls has table
\[
\begin{array}{c|ccc}
k&0&1&2\\\hline
\P(X=k)&1/10&6/10&3/10
\end{array}
\]
Thus $\E[X]=6/5$, $\E[X^2]=9/5$, and $\Var(X)=9/25$.
With replacement the count would instead be $\Bin(2,3/5)$, with the
same mean but variance $12/25$.
\end{example}

For the general moments, regard the sample as the first $m$ positions
in a uniform random ordering of the population.
Put $p=K/N$ and let $I_j$ indicate a marked object in position $j$.
Every position has marked probability $p$, so $\E[X]=mp$.
For $N\ge2$ and $i\ne j$,
\[
 \E[I_iI_j]=\frac{K(K-1)}{N(N-1)},\qquad
 \E[I_iI_j]-p^2=-\frac{p(1-p)}{N-1}.
\]
The variance-of-a-sum expansion from Chapter 4 gives
\[
 \Var(X)=mp(1-p)-m(m-1)\frac{p(1-p)}{N-1}
        =mp(1-p)\frac{N-m}{N-1}.
\]
For $1\le m\le N$, the last factor is at most one: sampling without
replacement has no more variance than the corresponding binomial count.
When $m=N$, the count is the fixed number $K$, so its variance is zero.
For a one-object population the count is constant as well.

\section{For later: geometric waiting times}
\label{sec:geometric}

Now allow a countably infinite list of possible values.
In independent trials with fixed success probability $0<p\le1$,
let $T$ count trials up to and including the first success. Then
\[
 \P(T=k)=(1-p)^{k-1}p,\qquad k=1,2,\ldots.
\]
For $0<p<1$, the first $k-1$ trials must fail and the $k$th succeed.
The geometric series sums the masses to one.
For $p=1$, $T=1$ with certainty. If $p=0$, there is no finite waiting
time to a success, so it is excluded from this distribution.

The event $T>r$ means $r$ initial failures, hence
$\P(T>r)=(1-p)^r$ for nonnegative integers $r$.
For $r=0$ the probability is one, including at $p=1$.
For $0<p<1$ and integers $s,t\ge0$,
\[
 \P(T>s+t\mid T>s)
 =\frac{(1-p)^{s+t}}{(1-p)^s}
 =(1-p)^t=\P(T>t).
\]
This is \emph{memorylessness}: after any number of failures, the
distribution of the additional wait is unchanged.
For $p=1$, only conditioning events of positive probability can be used.

\subsection*{Mean and variance from tail sums}
A positive integer $T$ satisfies
$T=\sum_{j\ge1}\ind_{\{T\ge j\}}$ and
$T^2=\sum_{j\ge1}(2j-1)\ind_{\{T\ge j\}}$.
First truncate these nonnegative sums; letting the truncation grow gives
\[
 \E[T]=\sum_{j\ge1}\P(T\ge j)=\sum_{j\ge1}(1-p)^{j-1}=\frac1p.
\]
The geometric-series calculation
$\sum_{j\ge1}j r^{j-1}=1/(1-r)^2$, reviewed in
Section~\ref{sec:calc-series}, similarly gives
\[
 \E[T^2]=\frac2{p^2}-\frac1p=\frac{2-p}{p^2},
 \qquad \Var(T)=\frac{1-p}{p^2}.
\]
These sums converge for $p>0$; at $p=1$ the values follow directly.
Waiting for a six on independent fair rolls gives mean $6$,
variance $30$, and $\P(T>5)=(5/6)^5$.

\begin{warningbox}
Some books count \emph{failures before} the first success.
That variable is $T-1$, supported on $0,1,\ldots$, with mean
$(1-p)/p$ and the same variance. State which convention you use.
\end{warningbox}

\section{For later: Poisson counts}
\label{sec:poisson}

For $\lambda>0$, the \emph{Poisson distribution} assigns
\[
 \P(X=k)=e^{-\lambda}\frac{\lambda^k}{k!},
 \qquad k=0,1,2,\ldots.
\]
The exponential series makes the total mass
$e^{-\lambda}\sum_{k\ge0}\lambda^k/k!=1$.
The limiting case $\lambda=0$ places all mass at zero.
Shifting indices in absolutely convergent sums gives
\[
 \E[X]=e^{-\lambda}\lambda\sum_{j\ge0}\frac{\lambda^j}{j!}
 =\lambda,\qquad
 \E[X(X-1)]=\lambda^2.
\]
Consequently $\E[X^2]=\lambda^2+\lambda$ and $\Var(X)=\lambda$.
Factorial moments make this calculation short, just as they did for hats.

The law arises from many independent rare opportunities. For fixed
$\lambda>0$ and fixed integer $k\ge0$,
\[
 \binom nk(\lambda/n)^k(1-\lambda/n)^{n-k}
 \longrightarrow e^{-\lambda}\frac{\lambda^k}{k!}.
\]
Indeed, $(n)_k/n^k\to1$ and $(1-\lambda/n)^n\to e^{-\lambda}$.
This explains a Poisson approximation to $\Bin(n,p)$ when $p$ is small
and $np$ is near $\lambda$; it does not make every count Poisson.
In particular, the model's equality of mean and variance is a consequence
to check, not a rule for all counts.

For $X$ Poisson with mean $3$, $\P(X\ge2)=1-4e^{-3}$.
The number of fixed points of a large uniform permutation approaches
Poisson with mean one, as derived in
Section~\ref{sec:calc-derangements}.
For a finite permutation this is an approximation; the exact law is
$\P(N=k)=\binom nkD_{n-k}/n!$.

\section{Coupon collection: waiting to see every type}
\label{sec:coupon-collector}

An album has $n\ge1$ types of coupon. Each purchase gives one coupon,
independently and uniformly chosen from the $n$ types; duplicates are
possible. Let $T_n$ be the number of purchases needed to see every type,
starting with an empty album. Equivalently, how many independent rolls
of a fair die are needed to see all six faces?

This reverses the birthday question. There we stopped at the first
repeat; here we keep going until nothing is missing.
With sampling without replacement, completion would take exactly $n$
draws. Repeated types are what make this problem interesting.

\subsection*{Break the wait into stages}
After $k$ distinct types have appeared, let $W_k$ count additional draws
up to and including the next new type, for $k=0,\ldots,n-1$.
There are $n-k$ unseen types. Until one appears, the chance of a new type
on each draw stays at
\[
 p_k=\frac{n-k}{n},\qquad
 W_k\sim\text{Geometric}(p_k),\qquad
 \E[W_k]=\frac{n}{n-k}.
\]
In particular $W_0=1$: the first draw is always new.
The collection time is the sum of these stages,
\[
 T_n=W_0+W_1+\cdots+W_{n-1}.
\]
Linearity of expectation now does the work. Define the \emph{harmonic
numbers} by $H_0=0$ and $H_n=\sum_{j=1}^n1/j$. Then
\begin{equation}
 \boxed{\E[T_n]=\frac nn+\frac n{n-1}+\cdots+\frac n1=nH_n.}
 \label{eq:coupon-mean}
\end{equation}
Every stage has finite expectation, so completion occurs with probability
one, although there is no fixed upper limit on the wait when $n\ge2$.
The last missing type alone takes $n$ further draws on average.

\begin{example}[Every face of a die]
For six equally likely types, the successive expected waits are
\[
 1,\quad\frac65,\quad\frac64,\quad\frac63,\quad\frac62,\quad6.
\]
Their sum is $6H_6=14.7$ rolls. Seeing the first five types takes $8.7$
rolls on average, and the last one takes another six.
The mean is an average over repetitions of the entire experiment,
not a deadline by which every collection will finish.
\end{example}

For a sense of scale, comparing rectangles with the area under $1/x$
gives, for $n\ge1$,
\[
 \log(n+1)=\int_1^{n+1}\frac{dx}{x}
 \le H_n\le1+\int_1^n\frac{dx}{x}=1+\log n.
\]
Thus $\E[T_n]$ grows like $n\log n$, rather than a fixed multiple of $n$.
Here and below, $\log$ is the natural logarithm. The integral comparison
can wait until Appendix~\ref{app:calculus}; the exact mean above uses
only geometric expectations and a finite sum.

\subsection*{Variance: justify independence before adding}
The expectation calculation did not require independence of the stages.
For the variance, we do need it. After any specified completed history,
future draws are still independent and uniform. The conditional law of
$W_k$ is always geometric with parameter $(n-k)/n$, regardless of which
types have appeared or how long earlier stages took. Applying the
multiplication rule to these conditional probabilities shows that
$W_0,\ldots,W_{n-1}$ are mutually independent.

The geometric variance formula therefore gives
\begin{align*}
 \Var(T_n)
 &=\sum_{k=0}^{n-1}\frac{1-p_k}{p_k^2}
   =\sum_{j=1}^n\frac{1-j/n}{(j/n)^2}\\
 &=\boxed{n^2\sum_{j=1}^n\frac1{j^2}-nH_n.}
\end{align*}
For a die, this is $38.99$ rolls squared, giving standard deviation
about $6.24$ rolls. For $n=1$, $T_1=1$ and the variance is zero.
For $n\ge2$, the entire collection time is not itself geometric:
the success probability changes each time a new type is found.

\subsection*{A deadline: count the types still missing}
Fix an integer $m\ge1$. Let $M_m$ be the number of types missing after
$m$ draws, and let $A_i$ be the event that type $i$ is missing.
Then
\[
 M_m=\sum_{i=1}^n\ind_{A_i},\qquad
 \P(A_i)=\left(1-\frac1n\right)^m,\qquad
 \E[M_m]=n\left(1-\frac1n\right)^m.
\]
These missing-type indicators are generally dependent. Linearity still
applies, and the union bound (equivalently, Markov at one) gives
\begin{equation}
 \P(T_n>m)=\P(M_m\ge1)
 \le\min\left\{1,n\left(1-\frac1n\right)^m\right\}
 \le\min\{1,ne^{-m/n}\}.
 \label{eq:coupon-tail-bound}
\end{equation}
We used $1-u\le e^{-u}$. Consequently, for $0<\delta<1$, any integer
$m\ge n(\log n+\log(1/\delta))$ guarantees completion probability
at least $1-\delta$. This is a sufficient bound, not the exact cutoff.

Inclusion--exclusion gives an exact answer: a specified set of $j$ types
is missing with probability $(1-j/n)^m$, so
\begin{equation}
 \P(T_n\le m)=\sum_{j=0}^n(-1)^j\binom nj
                         \left(1-\frac jn\right)^m,\qquad m\ge1.
 \label{eq:coupon-cdf}
\end{equation}
At $m=0$ the completion probability is zero and $M_0=n$.
For a die, the formula gives $\P(T_6\le15)\approx.6442$:
a deadline close to the mean is far from a guarantee.

\subsection*{Bonus: a Poisson number of draws}
There is also a direct connection to the Poisson distribution, entirely
within discrete probability. Choose $K\sim\Pois(t)$, where $t>0$,
independently of the coupon draws, and collect exactly $K$ coupons.
Let $C_i$ be the number of coupons of type $i$.
For nonnegative integers $c_1,\ldots,c_n$ with sum $s$, first condition
on $K=s$, then count the orderings of those coupon labels:
\begin{align*}
 \P(C_1=c_1,\ldots,C_n=c_n)
 &=\frac{e^{-t}t^s}{s!}\,
   \frac{s!}{c_1!\cdots c_n!}\left(\frac1n\right)^s\\
 &=\prod_{i=1}^n e^{-t/n}\frac{(t/n)^{c_i}}{c_i!}.
\end{align*}
The joint probability factors into Poisson probabilities. Thus the
$C_i$ are independent Poisson variables of mean $t/n$.
Each type is missing with probability $e^{-t/n}$, independently of
the other types. Hence
\[
 M_K\sim\Bin(n,e^{-t/n}),\qquad
 \P(T_n\le K)=\P(M_K=0)=\left(1-e^{-t/n}\right)^n.
\]
This trick is called \emph{Poissonization}. With a fixed number $m$ of
draws, the counts sum to $m$ and are dependent when $n\ge2$ and $m\ge1$.
Randomizing their total in this particular way makes them independent.
The product formula is exact for a Poisson number of draws; it is not
the fixed-$m$ answer in~\eqref{eq:coupon-cdf}.
No continuous waiting-time model is needed for this argument.

\section{Practice: choose and use a discrete model}
\label{sec:discrete-distribution-exercises}

\begin{exercisebox}
\begin{enumerate}
 \item For $X\sim\Bin(5,.3)$, find $\P(X\ge2)$, its mean, and variance.
 \exlevel{Fluency}
 \item Prove the geometric memoryless property, stating the allowed
 values and the condition that the conditioning event has positive
 probability. \exlevel{Reasoning}
 \item Approximate $\Bin(1000,.002)$ by a Poisson law and compare
 the exact and approximate $\P(X=0)$. \exlevel{Exploration}
 \item Choose a label uniformly from $1,\ldots,8$. Find its mean,
 variance, and probability of exceeding its mean. \exlevel{Fluency}
 \item An event has probability $.2$. Find the mean and variance
 of its indicator. Compare counting successes in ten independent
 repetitions with recording one indicator ten times. \exlevel{Reasoning}
 \item Draw two balls without replacement from three red and two blue
 balls. Derive the complete red-count distribution, mean, and variance.
 Repeat with replacement. \exlevel{Fluency}
 \item For each discrete family, describe one experiment that fits
 and one that violates a stated modeling assumption. \exlevel{Exploration}
 \item Roll a fair die independently until the first six. Find the
 mean and variance of the number of rolls and the probability it exceeds
 five. What is the mean number of failures? \exlevel{Fluency}
 \item For a Poisson variable of mean three, find $\P(X\ge2)$ and
 $\E[X(X-1)]$. Recover its variance. \exlevel{Fluency}
 \item A population has six marked objects among twenty. Take four
 draws. Compare the count's mean and variance with and without replacement.
 State why the probability models differ. \exlevel{Reasoning}
 \item Roll a fair die until all faces have appeared. Find the mean,
 variance, and expected additional wait when exactly two faces are still
 missing. \exlevel{Fluency}
 \item Let $T_{n,r}$ count draws until $r$ distinct types have appeared
 in the uniform coupon model. Show that
 $\E[T_{n,r}]=n(H_n-H_{n-r})$ for $1\le r\le n$.
 What is the expected remaining wait for a complete collection?
 \exlevel{Reasoning}
 \item For 100 equally likely coupon types, use the exponential bound
 in~\eqref{eq:coupon-tail-bound} to choose a whole number of draws
 guaranteeing at least $95\%$ completion probability. Explain why
 $100H_{100}$ is not such a guarantee. \exlevel{Reasoning}
 \item Independently choose $K\sim\Pois(12)$, then roll a fair die $K$ times.
 Find the distribution of the number of unseen faces and the probability
 of seeing every face. Compare with exactly 12 rolls. \exlevel{Exploration}
\end{enumerate}
\end{exercisebox}

\section{For later: continuous families}
\label{sec:continuous-families}

This section uses densities and integrals from
Section~\ref{sec:rv-continuous} and Appendix~\ref{app:calculus}.
It can be left until after the discrete expectation and variance material.

For $a<b$, Uniform$(a,b)$ assigns constant density $1/(b-a)$ to an interval. For $\lambda>0$, Exponential$(\lambda)$ has density $\lambda e^{-\lambda x}$ on $x\ge0$ and is the continuous memoryless law.

The normal density with mean $\mu$ and variance $\sigma^2>0$ is
\[
f(x)=\frac1{\sigma\sqrt{2\pi}}
\exp\!\left(-\frac{(x-\mu)^2}{2\sigma^2}\right).
\]
Standardization sends $X\sim\Normal(\mu,\sigma^2)$ to
\[
Z=\frac{X-\mu}{\sigma}\sim\Normal(0,1).
\]

\begin{center}
\begin{tikzpicture}[x=1.2cm,y=5cm]
  \draw[->,indigo] (-3.5,0)--(3.6,0) node[right] {$z$};
  \draw[mint,very thick,domain=-3.4:3.4,samples=80,smooth] plot (\x,{0.3989*exp(-\x*\x/2)});
  \fill[amber!35,domain=-1:1,samples=40] plot (\x,{0.3989*exp(-\x*\x/2)}) -- (1,0)--(-1,0)--cycle;
  \draw (-1,.012)--(-1,-.012) node[below] {$-1$};
  \draw (1,.012)--(1,-.012) node[below] {$1$};
  \node at (0,.19) {about $68\%$};
\end{tikzpicture}
\end{center}

For $\alpha,\beta>0$, the beta family on $[0,1]$ has density proportional to $x^{\alpha-1}(1-x)^{\beta-1}$. It is useful for proportions and becomes a natural prior for a Bernoulli probability.

\needspace{8\baselineskip}
\section{Choose by mechanism, not by silhouette}

A bell-shaped histogram does not prove normality. Ask what generated the
observations, what values are possible, whether trials are independent,
whether rates are stable, and what approximation error is acceptable.
A named distribution is useful because its assumptions describe an
experiment and its formulas then answer questions about that experiment.

\begin{exercisebox}[title={Later practice: continuous families}]
\begin{enumerate}[start=15]
 \item Find the median of an exponential distribution. \exlevel{Fluency}
 \item If $X\sim\Normal(10,9)$, express $\P(8<X<13)$ using the
 standard normal CDF. \exlevel{Fluency}
 \item Derive the mean and variance of Uniform$(a,b)$ for $a<b$.
 \exlevel{Reasoning}
\end{enumerate}
\end{exercisebox}
